% ID: FIS-DIN-NEW-002
% Titolo: Risultante di due forze oblique
% Disciplina: Fisica
% Area: Dinamica
% Argomento: Leggi di Newton
% Sottoargomento: Forze oblique e moto da fermo
% Classe: Terza liceo scientifico
% Difficolta: 4
% Tipo: problema_numerico
% Risultato: soluzione da aggiungere
% Tempo_stimato: 12 min
% Competenze: scomposizione vettoriale, secondo principio della dinamica, cinematica
% Prerequisiti: trigonometria, vettori, moto uniformemente accelerato
% Tag: fisica, dinamica, leggi_newton, forze_oblique, risultante
% Fonte: verifica_fisica_dinamica_anonimizzata
% Autore: Riccardo Carli
% Licenza: CC-BY-SA-4.0

\begin{esercizio}
Un punto materiale di massa \(m = 2{,}5\,\mathrm{kg}\) parte da fermo ed e
soggetto a due forze oblique. La prima forza ha modulo \(9{,}0\,\mathrm{N}\)
ed e diretta a \(30^\circ\) sopra l'asse \(+\hat x\); la seconda ha modulo
\(5{,}0\,\mathrm{N}\) ed e diretta a \(120^\circ\) rispetto a \(+\hat x\).

Calcola la risultante delle forze, esprimendola come
\[
\vec F_\text{tot} =
\begin{pmatrix}
F_x \\[4pt]
F_y
\end{pmatrix},
\]
e determina il vettore accelerazione. Sapendo che il moto inizia con velocita
nulla, trova lo spostamento del corpo dopo \(3{,}0\,\mathrm{s}\), espresso in
forma vettoriale e come modulo complessivo.
\end{esercizio}

\begin{soluzione}
% Soluzione da aggiungere.
\end{soluzione}

